\section{Introduction}

Retailers are already asking language models what to order. On InventoryBench, a
benchmark of 1{,}320 multi-period lost-sales replenishment instances, the
strongest published language-model strategy reaches a normalized reward of
$0.5380$ against $0.4447$ for the operations research policy it is asked to
improve~\cite{inventorybench}. That gap is why the approach is taken seriously.
It is also uninterpretable: the published strategy calls the model once per
period, lets it name an order quantity outright, and applies that quantity with
no test of whether the model was right, so three decisions are welded into one
design and the reported gain cannot be assigned to any of them.

Reproducing the baseline made this worse. Our arm~1 reimplements the published
capped base-stock policy and matches all 1{,}320 published order sequences bit
for bit, so it is the same policy rather than an approximation of one. Under the
benchmark's authoritative submission path that policy scores $0.5208$, not
$0.4447$, because the published figure came from the benchmark's legacy
simulator, which keeps a lost shipment visible in in-transit inventory forever
and so makes the baseline under-order after every loss. The motivating evidence
is therefore both bundled and partly artifactual, and the useful question is not
whether a language model helps on average but which of the three decisions does
the work.

\subsection*{The thesis}

A language model should not be asked what to order. It should be asked what
changed, once, in a vocabulary the system fixed in advance, and it should be
believed only after the data says so.

\collie{} implements that in three stages. A shared trigger decides \emph{when}:
an alert or a calibrated detector opens a proposal opportunity, a refractory
window suppresses repeats, and at most two proposals are permitted per episode. A
closed schema decides \emph{what}: at the proposal period $\tau_j$ the model
returns a \shockspec{}, a typed claim about the exogenous process drawn from a
registered vocabulary of families, directions, onset windows, magnitude bins and
durations, and it cannot write an order, a threshold, a likelihood, a falsifier,
or code. An anytime-valid \eprocess{} decides \emph{whether}: the hypothesis is
frozen at $\tau_j$ and influences no order until evidence arriving strictly after
$\tau_j$ drives the process above $1/\alpha_j$, and until then a deterministic
compiler is the only thing that ever touches the baseline controller's
parameters. Our notes name the design axis in one line: compile shocks, or always
inventory. Either the model's reading of an exception is compiled once into the
controller and then defended by prospective evidence, or the model is asked what
to order, every period, forever.

%% fig1_pipeline.tex -- Figure 1. The whole idea in one picture.
%%
%% FINAL FIGURE DESCRIPTION (for the illustrator, and for the record):
%% A left-to-right pipeline across the full text width, with a side panel.
%%   Left, stacked: the two trigger inputs, an operational alert and the demand
%%   telemetry stream. They meet in the trigger box, which is annotated with the
%%   refractory window of 5 periods and the cap of at most 2 proposals.
%%   The trigger opens a proposal opportunity at the stopping time tau_j.
%%   The language model is consulted exactly there and returns one typed
%%   ShockSpec, drawn with a lock icon to mark that it is frozen and immutable.
%%   The compiler maps the spec to one of 72 registered ControlConfig points.
%%   Between the compiler and the controller sits the e-process gate, drawn as a
%%   valve rather than an arrow, because that is the point of the figure: the
%%   compiled configuration is computed but does not flow until the evidence
%%   path E_{j,t}, accumulated only from observations after tau_j, crosses
%%   1/alpha_j. A dashed bypass path carries the baseline configuration straight
%%   to the controller and is the live path while the gate is closed.
%%   The controller emits q_t. A feedback arrow returns the realized observation
%%   Y_{t+1} to both the telemetry stream and the evidence path, which is what
%%   makes the closed-loop argument of Lemma 1 necessary rather than decorative.
%%   Right panel: the eleven-rung ladder as a vertical list, with arms 8, 9 and
%%   10 braced together and labeled "one physical call set, three activation
%%   policies", and arm 10 marked as the contribution.
%%
%% This file is the real figure, drawn in TikZ with base libraries only. It is
%% included by \input, not \includegraphics, so no external build step exists.

\begin{figure*}[!t]
\centering
\footnotesize
\begin{tikzpicture}[
  node distance=3.6mm,
  box/.style={draw, rounded corners=1pt, align=center, inner sep=2.2pt,
              minimum height=7.5mm, font=\footnotesize},
  src/.style={box, fill=black!4, minimum width=15mm},
  proc/.style={box, fill=black!6, minimum width=17mm},
  spec/.style={box, fill=black!10, minimum width=19mm, very thick},
  gate/.style={box, fill=black!14, minimum width=21mm, very thick},
  outbox/.style={box, fill=black!4, minimum width=11mm},
  rung/.style={align=left, font=\scriptsize, inner sep=1pt},
  flow/.style={-{Latex[length=1.6mm]}, thick},
  bypass/.style={-{Latex[length=1.6mm]}, thick, dashed},
  fb/.style={-{Latex[length=1.6mm]}, semithick, densely dotted},
  lbl/.style={font=\scriptsize, align=center, inner sep=1pt},
]

%% ---- inputs ---------------------------------------------------------------
\node[src] (alert) {operational\\alert};
\node[src, below=5mm of alert] (tele) {demand\\telemetry};

%% ---- trigger --------------------------------------------------------------
\node[proc, right=6mm of $(alert)!0.5!(tele)$] (trig) {trigger\\$\tau_j$};
\node[lbl, below=1.2mm of trig] {refractory 5\\$\le 2$ proposals};

%% ---- model ---------------------------------------------------------------
\node[proc, right=7mm of trig] (llm) {language\\model};

%% ---- the frozen spec ----------------------------------------------------
\node[spec, right=7mm of llm] (spec) {\shockspec{}\\\emph{frozen at} $\tau_j$};

%% ---- compiler ------------------------------------------------------------
\node[proc, right=7mm of spec] (comp) {compiler\\72 points};

%% ---- the gate ------------------------------------------------------------
\node[gate, right=7mm of comp] (gate)
  {\eprocess{} gate\\$E_{j,t} \ge 1/\alpha_j$};
\node[lbl, below=1.2mm of gate] {evidence from\\$\tau_j{+}1$ onward only};

%% ---- controller and order ------------------------------------------------
\node[proc, right=7mm of gate] (ctrl) {capped\\base-stock};
\node[outbox, right=6mm of ctrl] (order) {order $q_t$};

%% ---- flow ----------------------------------------------------------------
\draw[flow] (alert.east) -- (trig.west);
\draw[flow] (tele.east)  -- (trig.west);
\draw[flow] (trig)  -- (llm);
\draw[flow] (llm)   -- (spec);
\draw[flow] (spec)  -- (comp);
\draw[flow] (comp)  -- (gate);
\draw[flow] (gate)  -- (ctrl);
\draw[flow] (ctrl)  -- (order);

%% ---- the baseline bypass, live while the gate is closed ------------------
\coordinate (bypassL) at ($(comp.north) + (0,5.5mm)$);
\coordinate (bypassR) at ($(ctrl.north) + (0,5.5mm)$);
\draw[bypass] (comp.north) |- (bypassL) -- (bypassR) -| (ctrl.north);
\node[lbl, above=0.4mm] at ($(bypassL)!0.5!(bypassR)$)
  {baseline configuration, live until activation};

%% ---- closed loop ---------------------------------------------------------
\coordinate (fbY) at ($(tele.south) + (0,-8mm)$);
\draw[fb] (order.south) |- (fbY) -| (tele.south);
\node[lbl, above=0.4mm] at ($(fbY) + (26mm,0)$) {$Y_{t+1}$, and it also feeds the evidence path};

%% The ladder side panel was removed once Table~\ref{tab:ladder} carried the same
%% eleven rungs; keeping both duplicated the list and made this figure a third of
%% a page taller. The brace annotation that mattered is kept inline below.
\node[lbl, below=1.2mm of order] {arms 8, 9, 10 share\\one physical call set};

\end{tikzpicture}
\caption{The handshake. A shared trigger decides \emph{when} the model is
consulted, a closed schema decides \emph{what} it may say, and an
anytime-valid \eprocess{} decides \emph{whether} the answer touches an order.
The gate is drawn as a valve rather than an arrow because that is the point: the
compiled configuration is computed immediately and flows only after evidence
gathered strictly after $\tau_j$ drives $E_{j,t}$ above $1/\alpha_j$. The dashed
path is the baseline configuration, which is live until then. The dotted return
is what makes fact (F1) of Section~\ref{sec:verify} load-bearing rather than
decorative, since the controller's own orders shape the observations the test
consumes.}
\label{fig:pipeline}
\end{figure*}


One episode threads through the paper. A planner is told that a carrier has
suspended outbound scans on a lane and that containers booked over the last week
may not arrive at all. This is family~5 of our generator, a lost-shipment burst in
which orders dispatched in one to three consecutive periods are never delivered,
and it is a hard case on purpose: the demand stream is untouched, so the
calibrated detector fires at its null rate and the shortfall shows up in
telemetry only as an arrival that does not happen. Text is the only early signal,
and a wrong reading of it is expensive in both directions.

\subsection*{Contributions}

\begin{itemize}\itemsep0pt
\item \textbf{An audited testbed, and a divergence finding.} Our runner matches
      the official pipeline with a maximum absolute difference of exactly $0.0$
      over 648 comparisons, and arm~1 reproduces every published baseline order
      sequence bit for bit, which is how we established that the published score
      belongs to the legacy simulator (Section~\ref{sec:background}).
\item \textbf{\shockspec{} and its compiler.} A closed, frozen, typed vocabulary
      for what changed, plus a total deterministic map from a hypothesis to one
      point of a 72 configuration grid over a capped base-stock controller
      (Sections~\ref{sec:spec}, \ref{sec:compile}).
\item \textbf{Per-episode anytime-valid activation.} A discrete-mixture
      likelihood-ratio \eprocess{} started strictly after the proposal, with
      budget $\alpha_j = \alpha\,2^{-j}$ over at most two proposals, a
      family-by-family validity statement, and one explicitly prohibited shortcut
      (Section~\ref{sec:verify}).
\item \textbf{An eleven-rung ladder with honest cost accounting,} in which arms
      8, 9 and 10 share one physical set of proposal calls, so a difference
      between them cannot be a difference in calls (Section~\ref{sec:ladder}).
\end{itemize}

This draft is written against a running experiment, so every result cell is
marked \pending{} and inventoried, and Section~\ref{sec:notclaim} states the
single guarantee the system certifies and enumerates the readings of it that are
wrong.
